\chapter{Especificações e Requisitos}

O desenvolvimento conceitual e preliminar de uma aeronave comercial exige o alinhamento entre as demandas de mercado, traduzidas nos requisitos operacionais da missão, e as restrições impostas pelas autoridades certificadoras internacionais. Neste capítulo, serão apresentadas as especificações fundamentais do projeto, organizadas entre requisitos de missão, base normativa de aeronavegabilidade e a caracterização das missões típicas.

\section{Requisitos da Missão}
\label{section:Requisitos}

Conforme a especificação fornecida para o projeto da disciplina, a aeronave deve atender aos seguintes requisitos de desempenho e capacidade operacional:

\begin{itemize}
    \item \textbf{Missões de Projeto:} Transporte regional de passageiros (primária) e transporte de carga com o uso de contêineres e \textit{pallets} padronizados (ULDs) em configuração \textit{Quick Change} (secundária);
    
    
    \item \textbf{Capacidade de Passageiros:} 60 a 70 passageiros;
    
    \item \textbf{Peso Máximo de Decolagem (MTOW):} 50\,000 lb ($\approx 22\,680$ kgf);
    
    \item \textbf{Comprimento de Pista de Decolagem:} Máximo de 4\,200 ft ($\approx 1\,280$ m), ao nível do mar, condições ISA, em MTOW -- garantindo operacionalidade em aeroportos regionais de menor porte;
    
    \item \textbf{Comprimento de Pista de Pouso:} Máximo de 3\,000 ft ($\approx 914$ m), ao nível do mar, condições ISA, em MLW;
    
    \item \textbf{Velocidade Máxima de Cruzeiro:} 260 kt ($\approx 481$ km/h);
    
    \item \textbf{Alcance Máximo:} 800 NM ($\approx 1\,482$ km) com a quantidade determinada de passageiros e reservas IFR regulamentares (RBAC 121/FAR 121: alternado a 100 NM + 45 minutos de espera);
    
    \item \textbf{Teto Máximo Operacional (Certificado):} FL250 ($\approx 25\,000$ ft ou $7\,620$ m), em condições ISA -- altitude máxima certificada que otimiza a relação entre pressurização da cabine e arrasto de cruzeiro;
    
    \item \textbf{Configuração de Carga:} Compatibilidade com sistemas padronizados de carga aérea, incluindo contêineres e \textit{pallets} (ULDs), visando a conversão rápida (\textit{Quick Change}) e o aumento da taxa de utilização diária da frota;

    \item \textbf{Requisito de Conversão (QC):} O tempo de conversão da cabine de passageiros para carga (ou vice-versa) deve ser inferior a 60 minutos;

    \item \textbf{Ruído e Emissões:} Conformidade com ICAO Annex 16 (Chapter 14/Stage 4) e compatibilidade com 100\% SAF (\textit{Sustainable Aviation Fuel}).
\end{itemize}

\textbf{\textit{Nota:}} O teto de serviço do FL250 representa a altitude máxima de operação certificada da aeronave, o que não exige que todo voo de cruzeiro ocorra nessa altitude.

\section{Legislação e Requisitos de Aeronavegabilidade}

Com capacidade para até 70 passageiros e MTOW de 50\,000 lb, a aeronave enquadra-se na \textbf{Categoria Transporte}. Dessa forma, o projeto adota como base regulatória as normas de certificação emitidas pela ANAC, FAA e EASA, representadas respectivamente pelo RBAC n° 25, 14 CFR Part 25 e CS-25 \cite{rbac25, far25, cs25}.

Neste relatório, \emph{requisito de projeto} designa as especificações da Seção~\ref{section:Requisitos}, definidas para a missão, e \emph{norma} designa esses regulamentos de certificação. As seções da norma são citadas pelo número do 14 CFR Part 25 (\S 25.xxx), que segue a mesma numeração do RBAC 25 e da CS-25. Pesos são expressos em kgf (ou em lb, quando a correlação original usa unidades inglesas) e massas em kg.

Para o dimensionamento conceitual e verificação do projeto, destacam-se os seguintes requisitos normativos:

\begin{itemize}
    \item \textbf{Desempenho e Subida (\S 25.107 e \S 25.121):} Definição das velocidades operacionais e gradientes mínimos de subida em condição de motor inoperante;
    \item \textbf{Estabilidade e Controle Longitudinal (\S 25.143, \S 25.173 e \S 25.175):} Requisito de força no manche para variações de velocidade, assegurando autoridade do profundor e capacidade de trimagem via compensador ao longo de todo o envelope do Centro de Gravidade (CG);
    \item \textbf{Envelopes Estruturais (\S25.337 e \S25.341):} Fatores de carga limite entre $-1,0\,g$ e $+2,5\,g$ e tolerância a cargas de rajada;
    \item \textbf{Segurança de Cabine e Evacuação (\S25.803 e \S25.853):} Garantia de evacuação total da cabine em até 90 segundos com 50\% das saídas bloqueadas e conformidade dos materiais internos com os ensaios de inflamabilidade;
    \item \textbf{Ergonomia e Painel (\S25.771 e \S25.781):} Arranjo do cockpit no padrão Basic-T e visibilidade mínima para a tripulação.
    % ALTERADO: correção de digitação (de-acing/anti-acing -> de-icing/anti-icing)
    \item \textbf{Proteção contra Gelo e Condições Meteorológicas (\S25.1419):} A aeronave deve possuir sistemas de degelo/anticongelamento (\textit{de-icing/anti-icing}) certificados para operação sob condições conhecidas de formação de gelo.
\end{itemize}

Adicionalmente, a operação com cargas segue as especificações da IATA ULD Regulations \cite{iatauld} e as diretrizes da norma ISO 6870 \cite{iso4116} para padronização dimensional e travas de contenção. Recomendações práticas para distribuição de massa e amarração no porão/cabine são fundamentadas pela AC 120-85A da FAA \cite{ac12085}, enquanto as regras gerais de operação e navegação IFR atendem ao RBAC 91.

\section{Missões Típicas}

O projeto foi concebido para atender a duas demandas do mercado na aviação regional:

\begin{enumerate}
    \item \textbf{Transporte Regional de Passageiros:} Operação em rotas de 200 a 800 NM interligando cidades de médio porte à centros regionais. A operação é voltada para pistas curtas (até 4\,200 ft) e baixo tempo de \textit{turnaround}, acomodando de 60 a 70 passageiros com níveis adequados de conforto e pressurização até o FL250.

    \item \textbf{Transporte de Carga Aérea (\textit{Quick Change}):} Atendimento à logística regional e \textit{e-commerce}. A cabine adaptável permite a conversão para transporte de ULDs (\textit{Full Freighter}) ou operação mista (\textit{Kombi}), possibilitando o uso da frota em horários de baixa demanda de passageiros, como voos noturnos, para aumentar a utilização da aeronave. 
\end{enumerate}

Do ponto de vista operacional, a aeronave deverá ser capaz de realizar missões típicas dentro do alcance especificado de aproximadamente 800 milhas náuticas, permitindo cobrir uma ampla gama de rotas regionais.

Esse perfil de missão inclui etapas de subida, cruzeiro em altitudes médias e descida, com operação em aeroportos de menor porte, considerando as reservas de combustível regulamentares determinadas pelo RBAC 121/FAR 121 para operação sob Regras de Voo por Instrumentos (IFR) -- etapa de alternativa e 45 minutos de espera. Dessa forma, o projeto busca atender não apenas aos requisitos técnicos estabelecidos, mas também às necessidades reais do mercado, conciliando desempenho, flexibilidade e viabilidade operacional.